\documentclass[11pt]{article}
\usepackage[utf8]{inputenc}
\usepackage{amsmath,amssymb,booktabs,geometry,verbatim}
\usepackage{CJKutf8}
\geometry{margin=1in}

\title{CSP-S 2022 (Senior)}
\date{2022-10-29, 14:30--18:30 (4 hours)}

\begin{document}
\maketitle

\section*{Problem 1: Holiday Plan (\begin{CJK}{UTF8}{gbsn}假期计划\end{CJK})}

Input file \texttt{holiday.in}, output file \texttt{holiday.out}.
Time limit: 2 s. Memory limit: 512 MB.

Little Bear's map has $n$ points.
The point numbered $1$ is its home.
The points numbered $2, 3, \ldots, n$ are all scenic spots.
Some pairs of points are joined by a direct, bidirectional bus line.
If there are direct lines between $x$ and $z_1$,
between $z_1$ and $z_2$, \ldots,
between $z_{k-1}$ and $z_k$,
and between $z_k$ and $y$,
then we say that the trip between $x$ and $y$ is reachable with $k$ transfers.
In particular, if there is a direct line between $x$ and $y$,
we say that it is reachable with $0$ transfers.

The holiday is coming soon.
Little Bear plans to start from home, visit $4$ \emph{distinct} scenic spots,
and return home after completing $5$ legs:
home $\to$ spot A $\to$ spot B $\to$ spot C $\to$ spot D $\to$ home,
where each leg uses at most $k$ transfers.
There is no restriction on the points passed through while transferring:
they may be the home or scenic spots,
and the same point may be passed through repeatedly.
For example, on the leg spot A $\to$ spot B,
the points passed through while transferring may be the home, or spot C,
or points passed through while transferring on the leg spot D $\to$ home.

Suppose every scenic spot has a score.
Please help Little Bear plan a trip so that the sum of the scores of
the four \emph{distinct} scenic spots it visits is as large as possible.

\subsection*{Input}

Read the input from the file \texttt{holiday.in}.

The first line contains $3$ positive integers $n, m, k$,
which are the number of points on the map,
the number of pairs of points joined by a direct bidirectional line,
and the maximum number of transfers on each leg, respectively.

The second line contains $n-1$ positive integers,
the scores of the scenic spots numbered $2, 3, \ldots, n$.

Each of the next $m$ lines contains two positive integers $x, y$,
meaning that points $x$ and $y$ are directly joined by a road.
It is guaranteed that $1 \le x, y \le n$,
and that there are no multiple edges and no self-loops.

\subsection*{Output}

Write the output to the file \texttt{holiday.out}.

Output one positive integer:
the maximum possible sum of the scores of the $4$ distinct scenic spots that
Little Bear visits.

\subsection*{Sample 1 input}

\begin{verbatim}
8 8 1
9 7 1 8 2 3 6
1 2
2 3
3 4
4 5
5 6
6 7
7 8
8 1
\end{verbatim}

\subsection*{Sample 1 output}

\begin{verbatim}
27
\end{verbatim}

\subsection*{Explanation of sample 1}

When the planned trip is $1 \to 2 \to 3 \to 5 \to 7 \to 1$,
the sum of the scores of the $4$ scenic spots is $9+7+8+3 = 27$,
and it can be shown that this is the maximum.

The trip $1 \to 3 \to 5 \to 7 \to 8 \to 1$ has score sum $24$,
and the trip $1 \to 3 \to 2 \to 8 \to 7 \to 1$ has score sum $25$.
Both satisfy the requirements, but their score sums are not the largest.

The trip $1 \to 2 \to 3 \to 5 \to 8 \to 1$ has score sum $30$,
but the leg $5 \to 8$ needs at least $2$ transfers,
so it violates the requirement of at most $k = 1$ transfers.

The trip $1 \to 2 \to 3 \to 2 \to 3 \to 1$ has score sum $32$,
but it does not visit $4$ distinct scenic spots,
so it does not satisfy the requirements either.

\subsection*{Sample 2 input}

\begin{verbatim}
7 9 0
1 1 1 2 3 4
1 2
2 3
3 4
1 5
1 6
1 7
5 4
6 4
7 4
\end{verbatim}

\subsection*{Sample 2 output}

\begin{verbatim}
7
\end{verbatim}

\subsection*{Sample 3}

See \texttt{holiday/holiday3.in} and \texttt{holiday/holiday3.ans}
in the contestant's directory.

\subsection*{Constraints}

For all data, it is guaranteed that $5 \le n \le 2500$,
$1 \le m \le 10000$,
$0 \le k \le 100$,
and that the score of every scenic spot satisfies $1 \le s_i \le 10^{18}$.
It is guaranteed that at least one trip satisfying the requirements exists.

\begin{center}
\begin{tabular}{cccc}
\toprule
Test cases & $n \le$ & $m \le$ & $k \le$ \\
\midrule
$1 \sim 3$ & $10$ & $20$ & $0$ \\
$4 \sim 5$ & $10$ & $20$ & $5$ \\
$6 \sim 8$ & $20$ & $50$ & $100$ \\
$9 \sim 11$ & $300$ & $1000$ & $0$ \\
$12 \sim 14$ & $300$ & $1000$ & $100$ \\
$15 \sim 17$ & $2500$ & $10000$ & $0$ \\
$18 \sim 20$ & $2500$ & $10000$ & $100$ \\
\bottomrule
\end{tabular}
\end{center}

\section*{Problem 2: Strategy Game (\begin{CJK}{UTF8}{gbsn}策略游戏\end{CJK})}

Input file \texttt{game.in}, output file \texttt{game.out}.
Time limit: 1 s. Memory limit: 512 MB.

Little L and Little Q are playing a strategy game.

There is an array $A$ of length $n$ and an array $B$ of length $m$.
Based on them, a matrix $C$ of size $n \times m$ is defined,
satisfying $C_{ij} = A_i \times B_j$.
All indices start from $1$.

The game lasts $q$ rounds.
In each round, $4$ parameters $l_1, r_1, l_2, r_2$ are given in advance,
satisfying $1 \le l_1 \le r_1 \le n$ and $1 \le l_2 \le r_2 \le m$.

In the game, Little L first chooses an index $x$ between $l_1$ and $r_1$,
and then Little Q chooses an index $y$ between $l_2$ and $r_2$.
The score of the two players in this round is defined to be $C_{xy}$.

Little L's goal is to make this score as large as possible,
and Little Q's goal is to make this score as small as possible.
Both players are sufficiently clever and always play optimally.

Question: under the optimal strategies of the two players,
what is the score of each round?

\subsection*{Input}

Read the input from the file \texttt{game.in}.

The first line contains $3$ positive integers $n, m, q$,
which are the length of array $A$, the length of array $B$,
and the number of rounds, respectively.

The second line contains $n$ integers $A_i$, the elements of array $A$.

The third line contains $m$ integers $B_i$, the elements of array $B$.

Each of the next $q$ lines contains $4$ positive integers,
the values $l_1, r_1, l_2, r_2$ of that round.

\subsection*{Output}

Write the output to the file \texttt{game.out}.

Output $q$ lines, each containing one integer:
the score of Little L and Little Q in that round under optimal play.

\subsection*{Sample 1 input}

\begin{verbatim}
3 2 2
0 1 -2
-3 4
1 3 1 2
2 3 2 2
\end{verbatim}

\subsection*{Sample 1 output}

\begin{verbatim}
0
4
\end{verbatim}

\subsection*{Explanation of sample 1}

In this test, the matrix $C$ is as follows:

\begin{verbatim}
0 0
-3 4
6 -8
\end{verbatim}

In the first round, whether Little L chooses $x = 2$ or $x = 3$,
Little Q can choose some $y$ that makes the final score negative.
Therefore choosing $x = 1$ is optimal for Little L,
because then the score is certainly $0$.

In the second round, since Little L can choose $x = 2$,
Little Q can only choose $y = 2$, and the score is $4$.

\subsection*{Sample 2 input}

\begin{verbatim}
6 4 5
3 -1 -2 1 2 0
1 2 -1 -3
1 6 1 4
1 5 1 4
1 4 1 2
2 6 3 4
2 5 2 3
\end{verbatim}

\subsection*{Sample 2 output}

\begin{verbatim}
0
-2
3
2
-1
\end{verbatim}

\subsection*{Sample 3}

See \texttt{game/game3.in} and \texttt{game/game3.ans}
in the contestant's directory.

\subsection*{Sample 4}

See \texttt{game/game4.in} and \texttt{game/game4.ans}
in the contestant's directory.

\subsection*{Constraints}

For all data, $1 \le n, m, q \le 10^5$
and $-10^9 \le A_i, B_i \le 10^9$.
For every round, $1 \le l_1 \le r_1 \le n$
and $1 \le l_2 \le r_2 \le m$.

\begin{center}
\begin{tabular}{ccc}
\toprule
Test cases & $n, m, q \le$ & Special properties \\
\midrule
$1$ & $200$ & 1, 2 \\
$2$ & $200$ & 1 \\
$3$ & $200$ & 2 \\
$4 \sim 5$ & $200$ & none \\
$6$ & $1000$ & 1, 2 \\
$7 \sim 8$ & $1000$ & 1 \\
$9 \sim 10$ & $1000$ & 2 \\
$11 \sim 12$ & $1000$ & none \\
$13$ & $10^5$ & 1, 2 \\
$14 \sim 15$ & $10^5$ & 1 \\
$16 \sim 17$ & $10^5$ & 2 \\
$18 \sim 20$ & $10^5$ & none \\
\bottomrule
\end{tabular}
\end{center}

Here, special property 1 is: it is guaranteed that $A_i, B_i > 0$.

Special property 2 is: it is guaranteed that in every round,
either $l_1 = r_1$ or $l_2 = r_2$.

\section*{Problem 3: Star War (\begin{CJK}{UTF8}{gbsn}星战\end{CJK})}

Input file \texttt{galaxy.in}, output file \texttt{galaxy.out}.
Time limit: 2 s. Memory limit: 512 MB.

In this round of the interstellar war,
our side has built $n$ outposts in the universe,
joined by $m$ one-way wormholes.
We classify all wormholes whose endpoint is outpost $u$
as the wormholes of outpost $u$.

Amid the flames of war these wormholes can hardly last long;
enemy strikes may arrive at any time.
The effective strikes among them fall into two kinds:
\begin{enumerate}
\item The enemy destroys some wormhole.
This makes the two outposts it joins no longer directly reachable
through this wormhole,
but such a strike cannot destroy the two outposts it joins.
\item The enemy destroys some outpost.
Since the main technology of a wormhole is concentrated at its exit,
this causes all not-yet-destroyed wormholes of that outpost to be
destroyed together.
The wormholes starting from this outpost, however,
are \emph{not} destroyed.
\end{enumerate}
Note: destruction only makes a wormhole unusable;
it does not eliminate its existence.

To fight the enemy and maintain the links between the troops and the
outposts,
our side has developed two kinds of special forces responsible for
repairing wormholes:
\begin{itemize}
\item Type A special forces can repair one specific wormhole.
\item Type B special forces can repair all damaged wormholes of some
outpost.
\end{itemize}
Considering the nature of the enemy's strikes,
our side has not stockpiled too many strategic supplies at the outposts.
Therefore, as long as some wormhole of an outpost has been repaired and is
in a usable state, that outpost is also usable.

Our side has mastered a demanding spatial property.
Using this property, our warships can teleport along wormholes into the
enemy camp and carry out precision strikes.

To seize the best moment to launch a counteroffensive,
the headquarters must watch every change on the battlefield,
in order to find a moment at which a counteroffensive can be carried out.
The commander-in-chief holds that:
\begin{itemize}
\item If, starting from any one of our outposts,
with a suitably chosen route,
one can travel through wormholes infinitely many times
(the same outpost or the same wormhole may be passed through many times),
then this outpost \emph{can counterattack}.
\item To make the wormhole travel continuous and reduce as much as
possible the mass-energy loss of the warships when switching wormholes
at an outpost,
an outpost \emph{can travel continuously} if and only if
\emph{exactly one wormhole starting from that outpost is usable}.
\item If all of our outposts \emph{can counterattack}
and all of them \emph{can travel continuously},
then this moment is an excellent moment for a \emph{counteroffensive}.
\end{itemize}
The commander-in-chief gives you the order:
according to the information fed back from the battlefield in real time,
quickly tell him whether a \emph{counteroffensive} can be carried out at
the current moment.

\subsection*{Input}

Read the input from the file \texttt{galaxy.in}.

The first line of the input contains two positive integers $n, m$.

Each of the next $m$ lines contains two numbers $u, v$,
meaning a wormhole starting from outpost $u$ and going to outpost $v$.
It is guaranteed that $u \ne v$,
and that there are no two identical wormholes.
Initially all wormholes and outposts are intact.

The next line contains one positive integer $q$, the number of queries.

Each of the next $q$ lines describes a query or an operation.
First read a positive integer $t$, the type of the instruction:
\begin{itemize}
\item If $t = 1$, the next two integers $u, v$ mean that the enemy has
destroyed the wormhole from outpost $u$ to outpost $v$.
It is guaranteed that this wormhole exists and has not been destroyed.
\item If $t = 2$, the next integer $u$ means that the enemy has destroyed
outpost $u$.
If all wormholes of this outpost have already been destroyed,
this attack has no effect.
\item If $t = 3$, the next two integers $u, v$ mean that our side has
repaired the wormhole from outpost $u$ to outpost $v$.
It is guaranteed that this wormhole exists and is destroyed.
\item If $t = 4$, the next integer $u$ means that our side has repaired
outpost $u$.
If this outpost has no destroyed wormhole, this repair has no effect.
\end{itemize}
After each instruction is executed,
you need to decide whether a counteroffensive can be carried out.
If so, output \texttt{YES}, otherwise output \texttt{NO}.

\subsection*{Output}

Write the output to the file \texttt{galaxy.out}.

Output $q$ lines in total.
For each instruction, output whether a counteroffensive can be carried
out after this instruction is executed.

\subsection*{Sample 1 input}

\begin{verbatim}
3 6
2 3
2 1
1 2
1 3
3 1
3 2
11
1 3 2
1 2 3
1 1 3
1 1 2
3 1 3
3 3 2
2 3
1 3 1
3 1 3
4 2
1 3 2
\end{verbatim}

\subsection*{Sample 1 output}

\begin{verbatim}
NO
NO
YES
NO
YES
NO
NO
NO
YES
NO
NO
\end{verbatim}

\subsection*{Explanation of sample 1}

The states of the wormholes can be seen in the figure
``Figure 1: explanation of sample 1'' of the original statement,
in which the edges are the wormholes that exist and have not been
destroyed.
(The figure is not reproduced here.)

\subsection*{Sample 2}

See \texttt{galaxy/galaxy2.in} and \texttt{galaxy/galaxy2.ans}
in the contestant's directory.

\subsection*{Sample 3}

See \texttt{galaxy/galaxy3.in} and \texttt{galaxy/galaxy3.ans}
in the contestant's directory.

\subsection*{Sample 4}

See \texttt{galaxy/galaxy4.in} and \texttt{galaxy/galaxy4.ans}
in the contestant's directory.

\subsection*{Constraints}

For all data it is guaranteed that $1 \le n \le 5 \times 10^5$,
$1 \le m \le 5 \times 10^5$,
and $1 \le q \le 5 \times 10^5$.

\begin{center}
\begin{tabular}{ccccl}
\toprule
Test cases & $n$ & $m$ & $q$ & Special restrictions \\
\midrule
1, 2, 3 & $\le 10$ & $\le 20$ & $\le 50$ & none \\
4, 5, 6, 7, 8 & $\le 10^3$ & $\le 10^4$ & $\le 10^3$ & none \\
9, 10 & $\le 5 \times 10^5$ & $\le 5 \times 10^5$ & $\le 5 \times 10^5$
  & guaranteed no $t = 2$ and no $t = 4$ \\
11, 12 & $\le 5 \times 10^5$ & $\le 5 \times 10^5$ & $\le 5 \times 10^5$
  & guaranteed no $t = 4$ \\
13, 14, 15, 16 & $\le 10^5$ & $\le 5 \times 10^5$ & $\le 5 \times 10^5$
  & none \\
17, 18, 19, 20 & $\le 5 \times 10^5$ & $\le 5 \times 10^5$
  & $\le 5 \times 10^5$ & none \\
\bottomrule
\end{tabular}
\end{center}

\section*{Problem 4: Data Transmission (\begin{CJK}{UTF8}{gbsn}数据传输\end{CJK})}

Input file \texttt{transmit.in}, output file \texttt{transmit.out}.
Time limit: 3 s. Memory limit: 1024 MB.

Little C is designing the routing system of a computer network.

The network used for testing has $n$ hosts in total,
numbered $1 \sim n$ in order.
These $n$ hosts are joined by $n-1$ network cables;
the $i$-th cable joins hosts $a_i$ and $b_i$.
It is guaranteed that any two hosts are connected,
directly or indirectly, through finitely many cables.
Limited by the transmission power,
host $a$ can directly transmit information to host $b$ if and only if
the two hosts are connected, directly or indirectly,
through at most $k$ cables.

In a computer network, data transmission often needs several forwardings.
Suppose Little C needs to transmit data from host $a$ to host $b$
($a \ne b$).
Then it chooses some hosts used for the transmission,
$c_1 = a, c_2, \cdots, c_{m-1}, c_m = b$,
and forwards according to the following rule:
for all $1 \le i < m$, host $c_i$ sends the information directly to
$c_{i+1}$.

Every host needs some time to process information;
the $i$-th host needs $v_i$ units of time to process information.
Data travels through the network very quickly,
so the transmission time can be ignored.
Hence the time spent by the transmission process above is
$\sum_{i=1}^{m} v_{c_i}$.

Now there are $q$ data transmission requests in total;
the $i$-th request sends data from host $s_i$ to host $t_i$.
Little C wants to know, for each request,
the minimum number of units of time needed to complete the transmission.

\subsection*{Input}

Read the input from the file \texttt{transmit.in}.

The first line of the input contains three positive integers $n, Q, k$,
which are the number of hosts in the network, the number of requests,
and the transmission parameter, respectively.
The data guarantee $1 \le n \le 2 \times 10^5$,
$1 \le Q \le 2 \times 10^5$,
and $1 \le k \le 3$.

The second line of the input contains $n$ positive integers;
the $i$-th of them is $v_i$,
and it is guaranteed that $1 \le v_i \le 10^9$.

Each of the next $n-1$ lines contains two positive integers $a_i, b_i$,
meaning a cable joining hosts $a_i$ and $b_i$.
It is guaranteed that $1 \le a_i, b_i \le n$.

Each of the next $Q$ lines contains two positive integers $s_i, t_i$,
meaning a request to send data from host $s_i$ to host $t_i$.
It is guaranteed that $1 \le s_i, t_i \le n$ and $s_i \ne t_i$.

\subsection*{Output}

Write the output to the file \texttt{transmit.out}.

$Q$ lines, each containing one positive integer:
the minimum number of units of time needed by the $i$-th request during
transmission.

\subsection*{Sample 1 input}

\begin{verbatim}
7 3 3
1 2 3 4 5 6 7
1 2
1 3
2 4
2 5
3 6
3 7
4 7
5 6
1 2
\end{verbatim}

\subsection*{Sample 1 output}

\begin{verbatim}
12
12
3
\end{verbatim}

\subsection*{Explanation of sample 1}

For the first request, since hosts $4$ and $7$ need at least $4$ cables to
be connected,
the data cannot be transmitted directly between the two hosts,
and at least one forwarding is needed.
We let it be forwarded once at host $1$.
It is not hard to see that host $1$ needs only two cables to be connected
to each of hosts $4$ and $7$,
and the data processing time of host $1$ is only $1$,
the smallest among all hosts,
so the minimum transmission time is $4 + 1 + 7 = 12$.

For the third request, since hosts $1$ and $2$ need only $1$ cable to be
connected,
transmitting the data directly is optimal,
and the minimum transmission time is $1 + 2 = 3$.

\subsection*{Sample 2}

See \texttt{transmit/transmit2.in} and \texttt{transmit/transmit2.ans}
in the contestant's directory.

This sample satisfies the restrictions of test case 2.

\subsection*{Sample 3}

See \texttt{transmit/transmit3.in} and \texttt{transmit/transmit3.ans}
in the contestant's directory.

This sample satisfies the restrictions of test case 3.

\subsection*{Sample 4}

See \texttt{transmit/transmit4.in} and \texttt{transmit/transmit4.ans}
in the contestant's directory.

This sample satisfies the restrictions of test case 20.

\subsection*{Constraints}

For all test data,
$1 \le n \le 2 \times 10^5$,
$1 \le Q \le 2 \times 10^5$,
$1 \le k \le 3$,
$1 \le a_i, b_i \le n$,
$1 \le s_i, t_i \le n$,
and $s_i \ne t_i$.

\begin{center}
\begin{tabular}{ccccc}
\toprule
Test cases & $n$ & $Q$ & $k$ & Special property \\
\midrule
1 & $\le 10$ & $\le 10$ & $= 2$ & yes \\
2 & $\le 10$ & $\le 10$ & $= 3$ & yes \\
3 & $\le 200$ & $\le 200$ & $= 2$ & yes \\
4, 5 & $\le 200$ & $\le 200$ & $= 3$ & yes \\
6, 7 & $\le 2000$ & $\le 2000$ & $= 1$ & no \\
8, 9 & $\le 2000$ & $\le 2000$ & $= 2$ & no \\
10, 11 & $\le 2000$ & $\le 2000$ & $= 3$ & no \\
12, 13 & $\le 2 \times 10^5$ & $\le 2 \times 10^5$ & $= 1$ & no \\
14 & $\le 5 \times 10^4$ & $\le 5 \times 10^4$ & $= 2$ & yes \\
15, 16 & $\le 10^5$ & $\le 10^5$ & $= 2$ & no \\
17, 18, 19 & $\le 2 \times 10^5$ & $\le 2 \times 10^5$ & $= 2$ & no \\
20 & $\le 5 \times 10^4$ & $\le 5 \times 10^4$ & $= 3$ & yes \\
21, 22 & $\le 10^5$ & $\le 10^5$ & $= 3$ & no \\
23, 24, 25 & $\le 2 \times 10^5$ & $\le 2 \times 10^5$ & $= 3$ & no \\
\bottomrule
\end{tabular}
\end{center}

Special property: it is guaranteed that $a_i = i + 1$,
and $b_i$ is chosen uniformly at random from $1, 2, \ldots, i$.

\section*{Source}

The four problems of CSP-S 2022 (CCF CSP Senior, second round,
held 2022-10-29), as hosted on QOJ:
\begin{itemize}
\item Problem 1: \texttt{https://qoj.ac/problem/5025}
\item Problem 2: \texttt{https://qoj.ac/problem/5026}
\item Problem 3: \texttt{https://qoj.ac/problem/5027}
\item Problem 4: \texttt{https://qoj.ac/problem/5028}
\end{itemize}
The statements were translated from the official Chinese problem PDF
embedded in these pages.

\end{document}
